\section{Preparation limits and localization of least-favorable priors}
\label{sec:preparation-localization}
The contact structure has a second consequence, now for an arbitrary
number of training preparations. We separate a general localization
principle from the special geometry of the marked product family. An
exact finite-$N$ two-point prior is not assumed.

\begin{theorem}[Nearest-experiment limit and prior localization]
\label{thm:prior-localization}
Let $P_\theta$ be a continuous family on a finite alphabet of size $d$,
with compact metric parameter space, and let $Q$ be fixed. Put
\[
 d_Q(\theta)=\TV(Q,P_\theta),\qquad d_* =\min_\theta d_Q(\theta),
 \qquad C_* =\{\theta:d_Q(\theta)=d_*\}.
\]
For $N\ge1$, the unrestricted frozen-validation values satisfy
\begin{equation}\label{eq:nearest-limit}
 d_* -\sqrt{(d-1)/N}\le U_N(P,Q)\le d_*,\qquad
 U_N\uparrow d_*.
\end{equation}
Every least-favorable prior $\mu_N$ in Theorem~\ref{thm:dual} satisfies
\begin{equation}\label{eq:prior-excess}
 \int(d_Q-d_*)\,d\mu_N\le\sqrt{(d-1)/N}.
\end{equation}
Consequently every weak limit of such priors is supported on $C_*$.
For every $a>0$ their mass on $\{d_Q-d_*\ge a\}$ is at most
$\sqrt{(d-1)/N}/a$. More generally, an upper bound $b_N$ on twice the
expected total-variation error of a common empirical reconstruction
replaces $\sqrt{(d-1)/N}$ in both conclusions.
\end{theorem}
\begin{proof}
Let $\widehat P$ be the empirical law of the $N$ reports, and choose the
indicator $h_{\widehat P}$ maximizing $(Q-\widehat P)h$. If $h_\theta$
maximizes $(Q-P_\theta)h$, the two event inequalities give
\[
 (Q-P_\theta)h_{\widehat P}
 \ge d_Q(\theta)-2\TV(\widehat P,P_\theta).
\]
This inequality holds for the same parameter-independent response rule
at every parameter. For any law $P$ on $d$ atoms,
\[
 \mathbb E\TV(\widehat P,P)^2
 \le\frac{d}{4N}\left(1-\sum_xP(x)^2\right)
 \le\frac{d-1}{4N}.
\]
Thus its expected score is at least $d_Q(\theta)-\sqrt{(d-1)/N}$.
Every response has score at most $d_Q(\theta)$ at that parameter,
which proves the upper bound at a minimizing parameter. Ignoring the
last observation proves monotonicity and hence the limit.

Against a minimizing prior, the Bayes optimum is $U_N$. It is at least
the prior average of the same empirical rule. Therefore
\[
 U_N\ge \int d_Q\,d\mu_N-\sqrt{(d-1)/N}.
\]
Combine this with $U_N\le d_*$ to get \eqref{eq:prior-excess}.
Markov's inequality proves the mass estimate. The function $d_Q-d_*$
is nonnegative and continuous on the compact parameter space, so its
integral under any weak limit is zero. Such a limit is supported on
its zero set. The proof uses only the common reconstruction bound,
which proves the last assertion.
\end{proof}

The empirical estimate in \eqref{eq:nearest-limit} is the same
finite-alphabet mechanism underlying the retained decision-spectrum
classification. Its role here is different: it controls every
least-favorable prior, not merely a constructive selector. In particular,
finite-preparation saddle supports cannot remain far from the closest
candidate experiments as the preparation number grows.

\begin{theorem}[Intrinsic marked contact geometry]
\label{thm:marked-contact-geometry}
For every $\gamma\in I$, set $s_\gamma=\sqrt{(1+\gamma)/2}$ and
\[
 d_\gamma=2s_\gamma-1-\gamma/2,\qquad
 C_\gamma=\{(s_\gamma,s_\gamma),(1-s_\gamma,1-s_\gamma)\}.
\]
Then
\begin{equation}\label{eq:marked-distance}
 \min_{p,q}\TV(Q_\gamma,b(p,q))=d_\gamma,
 \qquad \mathop{\rm argmin}_{p,q}\TV(Q_\gamma,b(p,q))=C_\gamma.
\end{equation}
Moreover, for Euclidean distance in $[0,1]^2$,
\begin{equation}\label{eq:intrinsic-growth}
 \TV(Q_\gamma,b(p,q))-d_\gamma
 \ge\frac7{2300}\,\operatorname{dist}((p,q),C_\gamma)^2.
\end{equation}
The constant is uniform on $I$ and is derived from the experiment's
atom probabilities, not from a metric-entropy or small-ball assumption.
\end{theorem}
\begin{proof}
Put $a=(1-p)(1-q)$, $c=pq$, $l=(1-\gamma)/2$ and
$h=(1+\gamma)/2$. The common mass of target and candidate is
$\psi(a)+\psi(c)$, where
\[
 \psi(z)=
 \begin{cases}
 z,&0\le z\le l,\\
 z/2+(1-\gamma)/4,&l\le z\le h,\\
 1/2,&h\le z\le1.
 \end{cases}
\]
This follows by summing the smaller mass on each of the two marked
atoms corresponding to the report pair. Reflection
$(p,q)\mapsto(1-p,1-q)$ exchanges $a,c$; assume $c\ge a$.
If $c<1/4$, both terms are in the first segment of $\psi$, and their
sum is at most $1/2$.

For $c\ge1/4$, put $x=\sqrt c\ge1/2$. By the arithmetic-geometric
mean inequality,
\[
 a\le(1-x)^2\le1/4<l.
\]
Consequently the common mass is at most
$F(x)=(1-x)^2+\psi(x^2)$, with equality only when $p=q=x$.
On the three consecutive intervals with endpoints
$1/2,\sqrt l,\sqrt h,1$, this function is respectively
\[
 (1-x)^2+x^2,\quad
 (1-x)^2+x^2/2+(1-\gamma)/4,\quad
 (1-x)^2+1/2.
\]
The first two are strictly convex and the last strictly decreases.
The first increases away from $1/2$. Thus a maximum must occur at
$\sqrt l$ or $\sqrt h$. Their difference is
\begin{equation}\label{eq:overlap-endpoints}
 F(\sqrt h)-F(\sqrt l)
 =\gamma\left(\frac32-\frac2{\sqrt h+\sqrt l}\right)>\frac7{575}.
\end{equation}
Indeed $\sqrt h>39/50$, $\sqrt l>3/5$, and $\gamma\ge6/25$,
so the right side is greater than
$(6/25)(3/2-100/69)=7/575$.
The unique maximizing diagonal point on this side is $x=\sqrt h$.
Its common mass is $(1-\sqrt h)^2+1/2>1/2$, also excluding $c<1/4$.
Subtracting it from one gives $d_\gamma$. Reflection proves
\eqref{eq:marked-distance} and the assertion about all minimizers.

We record a uniform quantitative form. Write $s=\sqrt h$ and
$\Delta=7/575$. For $1/2\le x\le\sqrt l$, the common mass is at
most $F(\sqrt l)$, so its deficit from $F(s)$ is at least $\Delta$.
For $\sqrt l\le x\le s$, convexity puts $F$ below the chord between
its endpoint values. Since this interval has length at most one, its
deficit is at least $\Delta(s-x)$. For $s\le x\le1$ the deficit is
\[
 (1-s)^2-(1-x)^2=(x-s)(2-s-x)\ge(1-s)(x-s).
\]
Since $s<4/5$, this too is at least $\Delta|x-s|$.
Together the three intervals give $F(s)-F(x)\ge\Delta|x-s|$.

Let $z=(p-q)^2$. At fixed $c=x^2$ the additional loss from the
nondiagonal candidate is exactly
\[
 (1-x)^2-a=p+q-2x=\frac{z}{p+q+2x}\ge z/4.
\]
The geometric mean $x$ lies between $p$ and $q$, and hence
\[
 (p-s)^2+(q-s)^2\le4z+4(x-s)^2\le4z+4|x-s|.
\]
It follows that the total deficit is at least
$(\Delta/4)((p-s)^2+(q-s)^2)$, proving
\eqref{eq:intrinsic-growth} in this case. If $c<1/4$, the total variation is
at least $1/2$, whereas $d_\gamma<12/25$ from $s<4/5$ and
$\gamma\ge6/25$. The excess is therefore greater than $1/50$;
Euclidean squared distance in the square is at most two, and
$1/50>2(7/2300)$. This proves the remaining case. Reflection completes
the proof.
\end{proof}

\begin{corollary}[An all-preparation law and a limiting two-point prior]
\label{cor:marked-prior-limit}
For all $N\ge1$ and $\gamma\in I$,
\begin{equation}\label{eq:marked-all-N}
 d_\gamma-\sqrt{3/N}\le U_N(b,Q_\gamma)\le d_\gamma,
 \qquad U_N(b,Q_\gamma)\uparrow d_\gamma.
\end{equation}
Every least-favorable prior $\mu_{N,\gamma}$ satisfies
\begin{equation}\label{eq:marked-prior-moment}
 \int\operatorname{dist}(\theta,C_\gamma)^2d\mu_{N,\gamma}(\theta)
       \le\frac{2300}{7}\sqrt{3/N}.
\end{equation}
A least-favorable prior can be chosen invariant under simultaneous
complementation of the two parameters. Every such symmetric choice
satisfies, for the Euclidean quadratic transport distance,
\begin{equation}\label{eq:marked-prior-W2}
 W_2\left(\mu_{N,\gamma},
      \tfrac12\delta_{(s_\gamma,s_\gamma)}+
      \tfrac12\delta_{(1-s_\gamma,1-s_\gamma)}\right)
 \le\sqrt{2300/7}\,(3/N)^{1/4}.
\end{equation}
In particular, the limiting symmetric least-favorable law is unique.
\end{corollary}
\begin{proof}
Estimate the four-atom report-pair law from training and tensor it with
the known fair mark. Tensoring with a common probability preserves total
variation. The proof of Theorem~\ref{thm:prior-localization} therefore
uses $d=4$ rather than eight, proving \eqref{eq:marked-all-N} and an
excess bound $\sqrt{3/N}$. Combining the latter with
\eqref{eq:intrinsic-growth} gives \eqref{eq:marked-prior-moment}.

Simultaneously complement both report bits and the mark in every
training and validation word. This leaves $Q_\gamma$ invariant and
sends $b(p,q)$ to $b(1-p,1-q)$. The Bayes objective in
\eqref{eq:dual} is invariant under this transformation and convex in
the prior. Averaging a minimizing prior with its reflected image is
therefore again minimizing and is symmetric.

Map each parameter to a closest point of $C_\gamma$, choosing the two
points with equal probabilities at a tie. This is a Borel kernel,
equivariant under reflection. A symmetric prior is pushed forward to
the equal two-point law in \eqref{eq:marked-prior-W2}; the induced
coupling has squared transport cost exactly the left side of
\eqref{eq:marked-prior-moment}. Taking square roots proves the assertion.
\end{proof}

These conclusions identify the endpoint and limiting contact geometry
for every preparation number. They do not assert that the exact prior
has two atoms at every finite $N$, or that the displayed convergence
rates and constants are sharp. The finite-$N=1,2$ saddles are preserved
as exact results. The new theorem explains their relationship to an
intrinsically determined limiting contact set without extrapolating
their particular event tables to unproved values of $N$.
